%MB9T5
\loai{Ví dụ mẫu}
\begin{vidu}
Trong thí nghiệm Young về giao thoa ánh sáng, khoảng cách giữa hai khe là $2\ mm$, khoảng cách từ mặt phẳng chứa hai khe đến màn là $2\ m$. Ánh sáng đơn sắc dùng trong thí nghiệm có bước sóng $0,5 \ \mu m$. 
\begin{itemize}
\item [1.] Vùng giao thoa có bề rộng $5,97\ mm$. Trong vùng giao thoa số vân sáng là .......... và số vân tối là ..........
\item [2.] Hai điểm M và N trong vùng giao thoa, ở hai phía của vân trung tâm và cách trung tâm lần lượt $1,3\ mm$ và $2,9\ mm$. Số vân sáng trên MN là ..........
\item [3.] E và F thuộc vùng giao thoa có độ dài $EF=3\ mm$ và có 1 đầu là vân tối. Số vân sáng trên đoạn EF là ..........
\item [4.] P và Q thuộc vùng giao thoa có độ dài $PQ=4,7\ mm$ và có một đầu là vân sáng. Số vân tối trên đoạn PQ là ..........
\end{itemize}
\loigiai{
\begin{itemize}
\item [1.] Số vân sáng: $N_s=2\left[\dfrac{L}{2i}\right]+1=11$\\Số vân tối: $N_s=2\left[\dfrac{L}{2i}+0,5\right]=12$
\item [2.] Số vân sáng trên MN: $\dfrac{x_M}{i}\le k\le\dfrac{x_N}{i}\Leftrightarrow -2,6\le k\le 5,8$. Vậy có 8 vân sáng.  
\item [3.] $\dfrac{EF}{i}=6$ và EF có một đầu tối nên số vân sáng trên EF bằng 6
\item [4.] $\dfrac{PQ}{i}=9,4$ và PQ có một đầu sáng nên số vân tối bằng 9.
\end{itemize}
}
\end{vidu} \haidong{20}
\loai{Số vân sáng trong vùng giao thoa}
\begin{vidu}%[3P5-3.4K][Loai1]
Trong thí nghiệm Young về giao thoa ánh sáng, khoảng cách giữa hai khe là 0,5 mm, khoảng cách từ mặt phẳng chứa hai khe đến màn là 2 m. Ánh sáng đơn sắc dùng trong thí nghiệm có bước sóng $0,5 \ \mu m$. Vùng giao thoa trên màn rộng 26 mm. Số vân sáng là
\choice
{17}
{11}
{\True 13}
{15}
\loigiai{
\begin{itemize}
\item Khoảng vân: $i=\dfrac{\lambda D}{a}=\dfrac{0,5.10^{-6}.2}{0,5.10^{-3}}=2.10^{-3} \ m=2 \ mm.$
\item Số vân sáng: $N_s=1+2.\left[{\dfrac{L}{2i}}\right]=1+2.\left[{\dfrac{26}{2.2}}\right]=1+2.\left[{6,5}\right]=1+2.6=13$.
\end{itemize}
}
\end{vidu} \haidong{20}
\loai{Số vân tối trong vùng giao thoa}
\begin{vidu}%[3P5-3.4K][Loai2]
Một nguồn sáng đơn sắc S cách hai khe Young 0,2 mm phát ra một bức xạ đơn sắc có $\lambda =0,64\ \mu m$. Hai khe cách nhau $a=3\ mm$, màn cách hai khe 3 m. Miền vân giao thoa trên màn có bề rộng 12 mm. Số vân tối quan sát được trên màn là
\choice
{16}
{\True 18}
{19}
{17}
\loigiai{\begin{itemize}
\item Khoảng vân giao thoa: $i=\dfrac{\lambda D}{a}=\dfrac{0,64.3}{3}=0,64\ mm$.
\item Số vân tối quan sát được trên màn:
\begin{align*}
N_t=2.\left[\dfrac{L}{2i}+\dfrac{1}{2}\right]=2.\left[\dfrac{12}{2.0,64}+\dfrac{1}{2}\right]=2.\left[9,875\right]=2.9=18.
\end{align*}
\end{itemize}}
\end{vidu} \haidong{20}
\loai{Số vân sáng trên MN cho biết toạ độ}
\begin{vidu}%[3P5-3.4K][Loai4]
\nguon{QG - 2017} Trong thí nghiệm Young về giao thoa với ánh sáng đơn sắc có bước sóng $0,6\ \mu m$. Biết khoảng cách giữa hai khe là 0,6 mm, khoảng cách từ mặt phẳng chứa hai khe đến màn quan sát là 2 m. Trên màn, hai điểm M và N nằm khác phía so với vân sáng trung tâm, cách vân trung tâm lần lượt là 5,9 mm và 9,7 mm. Trong khoảng giữa M và N có số vân sáng là
\choice
{9}
{\True 7}
{6}
{8}
\loigiai{\begin{itemize}
\item Khoảng vân $i=\dfrac{D\lambda}{a}=2\ mm$.
\item Kết hợp với
$x_{M}< ki< x_{N}\Rightarrow -2,95< k< 4,85\Rightarrow$ có 7 giá trị.
\end{itemize}
}
\end{vidu} \haidong{20}
\loai{Số vân sáng - tối trên MN cho biết toạ độ}
\begin{vidu}%[3P5-3.3K][Loai5]
Trong thí nghiệm Young về giao thoa ánh sáng, các khe hẹp được chiếu sáng bởi ánh sáng đơn sắc. Khoảng vân trên màn là 1,2 mm. Trong khoảng giữa hai điểm M và N trên màn cùng một phía so với vân trung tâm, cách vân trung tâm lần lượt 2 mm và 4,5 mm, quan sát được
\choice
{\True 2 vân sáng, 2 vân tối}
{3 vân sáng, 2 vân tối}
{2 vân sáng, 3 vân tối}
{2 vân sáng, 1 vân tối}
\loigiai{
\begin{itemize}
\item Vị trí vân sáng trên màn: \bth{x_{s}=k\dfrac{\lambda.D}{a}=ki}.
\begin{itemize}
\item Vì M và N cùng phía so với vân trung tâm nên \bth{x_M>0;\ x_N>0}. 
\item Vì \bth{x_M< x_s< x_N} nên
\begin{align*}
x_M< ki< x_N\Rightarrow \dfrac{x_M}{i}< k< \dfrac{x_N}{i}\Leftrightarrow \dfrac{2}{1,2}< k< \dfrac{4,5}{1,2}\ (k\in \mathbb{Z})\Rightarrow k=2;3.
\end{align*}
\item Có 2 giá trị của k nên có hai vân sáng quan sát được trong khoảng MN.
\end{itemize}
\item Vị trí vân tối trên màn: \bth{x_t=(k+0,5)\dfrac{\lambda.D}{a}=(k+0,5)i}
\begin{itemize}
\item Vì M và N cùng phía so với vân trung tâm nên \bth{x_M>0;\ x_N>0}. 
\item Vì \bth{x_M< x_s< x_N} nên
\begin{align*}
x_M< (k+0,5)i< x_N\Rightarrow \dfrac{x_M}{i}< k+0,5< \dfrac{x_N}{i}\Rightarrow \dfrac{2}{1,2}-0,5< k< \dfrac{4,5}{1,2}-0,5\ (k\in \mathbb{Z})\Rightarrow k=2;3.
\end{align*}
\item Có 2 giá trị của k nên có hai vân tối quan sát được trong khoảng MN.
\end{itemize}
\end{itemize}
}
\end{vidu} \haidong{20}
\loai{Số vân sáng - 2 đầu sáng}
\begin{vidu}%[3P5-3.4K][Loai6]
Trên bề mặt rộng 7,2 mm của vùng giao thoa người ta đếm được 9 vân sáng (ở hai rìa là hai vân sáng). Tại vị trí cách vân trung tâm 14,4 mm là
\choice
{vân tối thứ 18}
{vân tối thứ 16}
{vân sáng bậc 18}
{\True vân sáng bậc 16}
\loigiai{
\begin{itemize}
\item Vì ở hai rìa là hai vân sáng nên khoảng vân
\begin{align*}
i=\dfrac{L}{n-1}=\dfrac{7,2}{9-1}=0,9\ mm.
\end{align*}
\item Tại M cách vân trung tâm 14,4\ mm ta có:$
\dfrac{x_M}{i}=\dfrac{14,4}{0,9}=16$ nên M là vân sáng bậc 16.
\item 
\end{itemize}
}
\end{vidu} \haidong{20}
\loai{Số vân sáng - 2 đầu tối}
\begin{vidu}%[3P5-3.4K][Loai7]
Trong thí nghiệm của Young về giao thoa ánh sáng đơn sắc, $a = 1 $ mm, $D =1 $ m. Biết hai điểm M, N trên màn quan sát cách nhau 3,6 mm có 6 vân sáng và tại M, N là vân tối. Bước sóng dùng trong thí nghiệm là
\choice
{\True 600 nm}
{500 nm}
{480 nm}
{560 nm}
\loigiai{
\begin{itemize}
\item Khoảng cách giữa 2 vân tối liên tiếp là 1i, giữa 1 vân tối và 1 vân sáng liên tiếp là 0,5i.
\item Vậy khoảng cách:
$MN=5i+0,5i+0,5i=6i=3,6\Rightarrow i=0,6\ mm\Rightarrow \lambda =\dfrac{ia}{D}=\dfrac{0,6. 1}{1}=0,6\ \mu m $
$=600\ nm $.
\end{itemize}
}
\end{vidu} \haidong{20}
\loai{Số vân sáng - 1 đầu tối - 1 đầu sáng}
\begin{vidu}%[3P5-3.2K][Loai8]
Trong thí nghiệm giao thoa ánh sáng Young người ta sử dụng ánh sáng đơn sắc. Giữa hai điểm M và N trên màn cách nhau 9\ mm chỉ có 5 vân sáng mà tại M là một trong 5 vân sáng đó, còn tại N là vị trí của vân tối. Xác định vị trí vân tối thứ 2 kể từ vân sáng trung tâm.
\choice
{\True $ \pm 3$ mm}
{$ \pm 0,3$ mm}
{$ \pm 0,5$ mm}
{$ \pm 5$ mm}
\loigiai{
\begin{itemize}
\item Ta có: $\Delta x=4i+0,5i\Rightarrow i=\dfrac{9}{4,5}=2\ mm\Rightarrow x_{t2}=\pm (2-0,5)i=\pm 3\ mm $.
\end{itemize}
}
\end{vidu} \haidong{20}

\loai{Số vân tối - sáng và ?}
\begin{vidu}%[3P5-3.4K][Loai11]
Trong thí nghiệm giao thoa Young với ánh sáng đơn sắc khoảng vân giao thoa là 0,5\ mm. Tại hai điểm M, N trên màn cách nhau 18,2 mm trong đó tại M là vị trí vân sáng. Số vân tối trên đoạn MN là
\choice
{\True 36}
{37}
{41}
{15}
\loigiai{
\begin{itemize}
\item Ta có: $N_t=\left[{\dfrac{MN-0,5i}{i}}\right]+1=\left[{\dfrac{18,2}{0,5}-0,5}\right]+1=36$.
\end{itemize}
}
\end{vidu} \haidong{20}
\loai{Tổng hợp}
\begin{vidu}%[3P5-3.4K][Loai12]
Trong thí nghiệm giao thoa Young, trên màn quan sát hai vân sáng đi qua hai điểm M và P. Biết đoạn MP dài 7,2\ mm đồng thời vuông góc với vân trung tâm và số vân sáng trên đoạn MP nằm trong khoảng từ 11 đến 15. Tại điểm N thuộc MP, cách M một đoạn 2,7\ mm là vị trí của một vân tối. Số vân tối quan sát được trên MP là
\choice
{11}
{\True 12}
{13}
{14}
\loigiai{
\begin{itemize}
\item Số vân sáng trên đoạn MP:
$11<N_{MP}=\dfrac{MP}{i}+1<15\Rightarrow 0,514\ mm<i<0,72\ mm$.
\item Vì M là vân sáng và N là vân tối nên:
\begin{align*}
MN=\left({n+0,5}\right)i
\Rightarrow 2,7=\left({n+0,5}\right)i\Rightarrow i=\dfrac{2,7}{n+0,5}\xrightarrow{{0,514<i<0,72}}3,25<n<4,75\Rightarrow n=4
\Rightarrow i=\dfrac{2,7}{4+0,5}=0,6\ mm.
\end{align*}
\item Số vân tối trên đoạn MP: $N_t=\dfrac{MP}{i}=\dfrac{7,2}{0,6}=12$.
\end{itemize}
}
\end{vidu} \haidong{20}


